\begin{rubric}{工作经历}
\entry*[2024年8月 -- 至今]%
	\textbf{波士顿学院（Boston College，U.S. News 全美排名 36）计算机系}\hspace{3em}博士后研究员
 \begin{itemize}
     \item 合作导师：\href{https://donglaiw.github.io/}{魏东来} 教授。
     \item 研究方向：面向线粒体、细胞核及多细胞器的三维体电镜实例分割与结构解析。
     \item 主要工作与成果：
     \begin{itemize}
         \item 开展 EM 三维线粒体、细胞核及全细胞器实例分割的方法研究与算法设计。
         \item 牵头推进 MitoEM 2.0 基准构建与论文撰写，形成数据集、算法与评测相结合的研究闭环。
         \item 在 CellMap Segmentation Challenge（eFIB-SEM 多细胞器分割，40+ 类别）中获得第一名。
         \item 正在推进全细胞器实例分割、骨架引导结构学习、多分辨率专家混合模型等工作，拟投稿 \textit{IEEE Transactions on Medical Imaging}。
     \end{itemize}
 \end{itemize}
\entry*[2017年7月 -- 2019年7月]%
	\textbf{美团点评}\hspace{5em}搜索与推荐算法工程师
 \begin{itemize}
     \item 主要工作：
     \begin{itemize}
         \item 负责旅游业务排序模型研究与迭代。
         \item 面向点击率（CTR）与转化率（CVR）提升，开展学习排序与召回优化。
         \item 相关技术包括 Graph Embedding、DNN、Wide\&Deep 等。
     \end{itemize}
 \end{itemize}

 % \entry*[May 2017 -- Jul 2017]%
	% \textbf{Sinovo-tech }\hspace{8em}Computer Vision Algorithm Engineer.
 % \begin{itemize}
 %     \item Participate in the development of video understanding algorithm in intelligent transportation field.
 %     \item \textbf{Related algorithm}: Faster RCNN, YoLo, SSD
 % \end{itemize}
%
% Blank lines result in extra space!
%
%
\end{rubric}
